\begin{lemma}[Meromorphic quotient as divisor germs]
\label{lemma-meromorphic-quotient-divisor-germs}
\begin{reference}
\cite[Chapter III, \S 3.4, p. 160]{Tohoku}
\end{reference}
Let $X$ be a Noetherian integral scheme and let $\mathcal{K}_X$ be the
sheaf of meromorphic functions. Taking the quotient in the category of
sheaves of abelian groups, the stalk at $x \in X$ is
\[
(\mathcal{K}_X^*/\mathcal{O}_X^*)_x
\cong K^*/\mathcal{O}_{X,x}^*,
\]
where $K$ is the function field of $X$. We call this quotient the sheaf of
germs of locally principal divisors: a germ is represented by a nonzero
meromorphic function, with two representatives identified when their ratio
is a unit in the local ring. If, in addition, every local ring of $X$ is a
unique factorization domain and $\mathcal{D}_X$ is the sheaf of Weil
divisors, then the divisor map induces a canonical isomorphism
\begin{equation}
\label{equation-meromorphic-quotient-divisor-germs}
\mathcal{K}_X^*/\mathcal{O}_X^*
\xrightarrow{\ \sim\ }
\mathcal{D}_X.
\end{equation}
For every $x \in X$, with $A = \mathcal{O}_{X,x}$ and $K$ its fraction
field, the induced stalk isomorphism is
\[
(\mathcal{K}_X^*/\mathcal{O}_X^*)_x
\cong K^*/A^*
\cong
\bigoplus_{\substack{\mathfrak p \subset A\\
\operatorname{height}(\mathfrak p)=1}}
\mathbf Z[\mathfrak p].
\]
Thus under the locally factorial hypotheses the locally principal divisor
germs are exactly the Weil-divisor germs.
\end{lemma}

\begin{proof}
The meromorphic-functions lemma identifies the stalk of $\mathcal{K}_X$ at
every point of the integral scheme with the function field $K$. Hence the
stalk of the sheaf quotient is $K^*/A^*$, which is exactly the stated
locally principal divisor-germ relation. Under the additional locally UFD
hypothesis, the exact sequence of Lemma
\ref{lemma-locally-UFD-divisor-resolution} identifies its cokernel with
the quotient in (\ref{equation-meromorphic-quotient-divisor-germs}).
Taking stalks and using the computation in that lemma then gives
\[
(\mathcal{K}_X^*/\mathcal{O}_X^*)_x \cong K^*/A^*
\cong
\bigoplus_{\substack{\mathfrak p \subset A\\
\operatorname{height}(\mathfrak p)=1}}
\mathbf Z[\mathfrak p].
\]
The last group is the stalk of $\mathcal{D}_X$. Unique factorization lets
one choose a meromorphic representative for every finite divisor germ, so
the stalkwise description is exactly the stated locally principal
interpretation.
\end{proof}

\begin{remark}[Algebraic scope of the quotient identification]
\label{remark-meromorphic-quotient-algebraic-boundary}
The preceding identification is conditional on the displayed algebraic
hypotheses and uses multiplicative sheaves of abelian groups. It does not
assert an analytic quotient, a GAGA comparison, or any descent statement;
those require separate analytic input.
\end{remark}
